\section{Introduction}

Bitcoin is governed by nobody and changed by somebody. Its consensus rules are enforced by every node that runs the software, so no person or body can alter them by decree; yet the rules do change---P2SH in 2012, CHECKLOCKTIMEVERIFY in 2015, Segregated Witness in 2017, Taproot in 2021---and each change was proposed, argued, revised, activated or abandoned by identifiable people writing on a public mailing list. The Bitcoin Improvement Proposal (BIP) process \cite{bip2} is the paper trail of those changes: a BIP is a design document that describes a change to the protocol, the peer-to-peer network, wallets or the process itself; it is discussed on the bitcoin-dev list, assigned a number by a BIP editor, and moves through a small set of statuses (Draft, Proposed, Final, Active, Deferred, Rejected, Withdrawn, Replaced, Obsolete; since the 2025 revision of the process, Draft, Deployed, Complete and Closed \cite{bip2}). Publication of a BIP does not mean that the community has accepted it \cite{bitcoinbips}; whether a proposal is adopted depends on discussion, on implementations, on the choices of miners, node operators, wallets and exchanges, and---for consensus changes---on an activation mechanism that has itself been the subject of some of the community's bitterest disputes \cite{bip8, bip148}.

This makes Bitcoin an unusually good place to study influence. In a project such as Python, a decision has a formal owner: the BDFL, a delegate, or the Steering Council can pronounce a proposal accepted \cite{pep1, pep13}. Bitcoin has no such office. It has maintainers of the reference implementation with merge access, a lead maintainer until January 2022, editors who administer the BIP repository, proposal authors, companies that employ most of the active developers, miners whose signalling activates soft forks, and a large community of node operators, wallet developers and users who can refuse to run new code. Prior work has characterised this arrangement as ``invisible politics'' \cite{defilippi2016}, as ``senatorial'' governance by competing bureaucratic parties \cite{parkin2019}, as a fiduciary relationship between developers and users \cite{walch2019}, and, quantitatively, as a system in which a handful of people account for most of the discussion \cite{azouvi2018}. What such accounts leave open is \emph{how} influence is exercised in the discussions themselves: with which arguments, by whom, at what point of a proposal's life, and through which of the contested forms---decisions by fiat, corporate stakes, gatekeeping, threats to fork, incivility, off-list agreements, procedural control---that the block-size war of 2015--2017 made visible to everyone \cite{bier2021}.

This paper applies Influence Miner, the tool and taxonomy introduced for Python \cite{influenceMinerPython}, to the complete public record of Bitcoin's proposal discussions: every message on the bitcoin-dev mailing list from its creation in June 2011 to September 2026, linked to the BIPs it discusses, together with the status history of every BIP from the \texttt{bitcoin/bips} repository and role rosters derived from the \texttt{bitcoin/bitcoin} repository. Influence Miner classifies sentence-level influence signals into thirteen base mechanisms (strategic, operational, functional, tactical, authority, compatibility, security, standards, ecosystem, economic, organizational, coalition and user demand) and seven controversial ones (unilateral, corporate interest, gatekeeping, exit threat, incivility, backchannel and procedural control), and attaches to each signal a direction, a scope, a target, the author's role, the proposal's outcome and the signal's timing relative to the decision. The taxonomy was designed with Bitcoin in mind---its cue lists already carried soft forks, miners, wallets, exchanges, fee markets and consensus failures---and this paper is its first test on the community those cues were written for.

The paper makes four contributions. First, it delivers the Bitcoin corpus itself: a reproducible pipeline that turns the public-inbox archive of bitcoin-dev, the BIP repository's history and the reference implementation's commit history into the same relational form that the DeMaP Miner family of tools uses for Python, with 24,391 messages, 205 BIPs and role labels for every message. Second, it reports the first sentence-level influence profile of Bitcoin's proposal discussions and answers six research questions about mechanisms, actors, eras, outcomes, the overlap with preference signals and controversial influence. Third, it compares Bitcoin with Python on the same instrument and finds the differences the taxonomy predicted---more ecosystem, economic and security influence, weaker internal authority, more exit and coalition pressure---together with differences it did not predict. Fourth, it uses the tool to nominate the landmark influence acts of each year of Bitcoin's history and sets them against the events the community itself remembers.

As with the Python study, the labels are produced by transparent cue rules that have not yet been validated against a gold standard; the results are a baseline and a demonstration that the pipeline transfers, not final measurements of influence in Bitcoin.

